%mac
\loai{Ví dụ mẫu}
\begin{vidu}%[1P1-7.3K][Loai1]
Một vật được thả từ độ cao $h = 80\ m$. Bỏ qua sức cản không khí. Lấy $g = 10\ m/s^2$. 
\begin{vdc}
Thời gian vật rơi cho đến khi chạm đất là bao nhiêu giây?
\shortans[oly]{4}
\loigiai{
Thời gian vật rơi cho đến khi chạm đất: $t=\sqrt{\dfrac{2h}{g}}=\sqrt{\dfrac{2.80}{10}}=4\ s$
}
\end{vdc}
\begin{vdc}
Vận tốc của vật khi vừa chạm đất có độ lớn là bao nhiêu $m/s$?
\shortans[oly]{40}
\loigiai{
Vận tốc của vật khi vừa chạm đất (vận tốc của vật lúc t = 4 s): 	$v=gt=10.4=40\ m/s$
}
\end{vdc}
\begin{vdc}
Độ cao của vật sau khi rơi $2\ s$ là bao nhiêu mét?
\shortans[oly]{60}
\loigiai{
Độ cao của vật sau khi rơi 2s: $h'=h-S=h-\dfrac{1}{2}gt^2=60\ m$
}
\end{vdc}
\end{vidu}



\begin{vidu}%[1P1-7.3K][Loai1]
Một vật rơi tự do. Trong $4\ s$ cuối cùng rơi được $320\ m$. Lấy $g = 10\ m/s^2$.
\begin{vdc}
Thời gian rơi là bao nhiêu giây?
\shortans[oly]{10}
\loigiai{
\begin{itemize}
\item Gọi t là thời gian vật rơi trong toàn bộ quãng đường h: $h=\dfrac{1}{2}gt^2=5t^2$
\item Quãng đường vật rơi được trong khoảng thời gian ($t - 4$) giây là $h_1$:\\
$h_1=\dfrac{1}{2}g{\left({t-4}\right)}^2=5{\left({t-4}\right)}^2=5t^2-40t+80$
\item Theo đề ra ta có: $h-h_1=320\Rightarrow 5t^2-\left({5t^2-40t+80}\right)=320\Rightarrow t=10\ s$
\end{itemize}
}
\end{vdc}
\begin{vdc}
Vận tốc của vật ngay trước khi chạm đất có độ lớn là bao nhiêu $m/s$?
\shortans[oly]{100}
\loigiai{
\begin{itemize}
\item Vận tốc khi chạm đất: $v=gt=10.10=100\ m/s$
\end{itemize}
}
\end{vdc}
\end{vidu}
\loai{Thời gian vật rơi mét thứ n}
\begin{vidu}%[1P1-7.3K][Loai1]
Một vật rơi tự do tại nơi có $g = 10\ m/s^2$, thời gian rơi là $20\ s$.
\begin{vdc}
Quãng đường rơi trong $1\ s$ đầu là bao nhiêu mét?
\shortans[oly]{5}
\loigiai{
\begin{itemize}
\item Quãng đường của vật rơi tự do: $s=\dfrac{1}{2}gt^2=5t^2$
\item Quãng đường vật rơi được trong $t = 20$ s: $s=\dfrac{1}{2}gt^2=5t^2=5.20^2=2000\ m$
\item Quãng đường rơi được trong thời gian 1 s đầu: $s_1=\dfrac{1}{2}gt^2=5t^2=5.1^2=5\ m$
\end{itemize}
}
\end{vdc}

\begin{vdc}
Quãng đường rơi trong thời gian $1\ s$ cuối là bao nhiêu mét?
\shortans[oly]{195}
\loigiai{
\begin{itemize}
\item Quãng đường vật rơi được trong thời gian 19 s đầu: $s_2=\dfrac{1}{2}gt^2=5t^2=5.19^2=1805\ m$
\item Quãng đường rơi được trong 1 s cuối: $s-s_1=2000-1805=195\ m$
\end{itemize}
}
\end{vdc}

\begin{vdc}
Thời gian rơi trong $1\ m$ đầu là bao nhiêu mili giây (làm tròn kết quả đến chữ số hàng đơn vị)?
\shortans[oly]{447}
\loigiai{
\begin{itemize}
\item Thời gian rơi trong 1 m đầu: $t=\sqrt{\dfrac{2h}{g}}=\sqrt{\dfrac{2.1}{10}}=0,447\ s$
\end{itemize}}
\end{vdc}

\begin{vdc}
Thời gian rơi trong $1\ m$ cuối là bao nhiêu mili giây (làm tròn kết quả đến chữ số hàng đơn vị)?
\shortans[oly]{5}
\loigiai{
\begin{itemize}
\item Thời gian rơi trong (2000 - 1) m đầu: $t_{1999}=\sqrt{\dfrac{2h}{g}}=\sqrt{\dfrac{2.1999}{10}}=19,995\ s$
\item Thời gian rơi trong 1m cuối: $\Delta t=t_{2000}-t_{1999}=20-19,995=0,005\ s$
\end{itemize}}
\end{vdc}
\end{vidu}



\loai{Xác định các đại lượng cơ bản của rơi tự do}
\begin{vidu} %[1P1B7-2][Loai1]
Một vật rơi tự do từ độ cao $19,6\ m$ xuống đất. Lấy $g = 9,8\  m/s^2$. Thời gian để vật đi hết đoạn đường đó là
\choice
{$t = 4\ s$}
{\True $t = 2\ s$}
{$t = 6\ s$}
{$t = 8\ s$}
\haidong{2}
\loigiai{
Ta có: $t = \sqrt{\dfrac{2S}{g}} = \sqrt{\dfrac{2.19,6}{9,8}} = 2\ s$.
}
\haidong{7} \end{vidu}
\begin{vidu} %[1P1B7-2][Loai1] 
Một vật rơi tự do từ độ cao $5\ m$. Cho $g = 9,8\  m/s^2$. Tốc độ của vật khi chạm đất là
\choice
{$1,01\  m/s$}
{$1,98\  m/s$}
{$7,0\  m/s$}
{\True $9,9\  m/s$}
\haidong{2}
\loigiai{	      
Vận tốc của vật khi chạm đất là: $v_t=\sqrt {2gh}=\sqrt{2.9,8.5}\approx 9,9\  m/s$.
}
\haidong{7} \end{vidu}
\loai{Xác định các đại lượng cơ bản của rơi tự do - bài toán ngược}
\begin{vidu}%[1P1-7.2B][Loai1]
Một hòn đá được thả từ miệng một giếng cạn đến đáy giếng mất $3\ s$. Bỏ qua sức cản không khí. Lấy $g = 9,8\ m/s^2$. Độ sâu của giếng là
\choice
{$90,3\ m$}
{$45,8\ m$}
{$30,4\ m$}
{\True $44,1\  m$}
\haidong{4}
\loigiai{
\begin{itemize}
\item Chọn gốc tọa độ tại miệng giếng, chiều dương của hệ trục hướng xuống dưới.
\item Ta có: $y=\dfrac{gt^2}{2}=4,9t^2\ m$.
\item Độ sâu của giếng là: $y=4,9.3^2=44,1\ m$.
\end{itemize}	
}
\haidong{7} \end{vidu}
\begin{vidu}%[1P1-7.2B][Loai1]
Thời gian rơi của một vật được thả từ trên cao xuống đất là $4\ s$. Bỏ qua sức cản không khí. Lấy $g = 10\ m/s^2$. Tốc độ lúc chạm đất của vật là
\choice
{$10\ m/s$}
{$20\ m/s$}
{\True $40\ m/s$}
{$80\ m/s$}
\haidong{4}
\loigiai{
\begin{itemize}
\item Chọn gốc tọa độ tại vị trí vật rơi, chiều dương của trục tọa độ hướng xuống.
\item Ta có: $v=gt=10t \ m/s$.
\item Vận tốc của vật ngay trước khi chạm đất là: $v = 10.4 = 40$ m/s.
\end{itemize}	
}
\haidong{7} \end{vidu}
\begin{vidu}%[1P1-7.2B][Loai1]
Một người thả vật rơi một vật chạm đất có $v = 50\ m/s$. Bỏ qua sức cản không khí. Lấy $g = 10\ m/s^2$. Độ cao của vật sau khi đi được $3\ s$ là
\choice
{$45\ m$}
{\True $80\ m$}
{$100\ m$}
{$125\ m$}
\haidong{4}
\loigiai{
\begin{itemize}
\item Chọn gốc tọa độ tại vị trí vật bắt đầu rơi, chiều dương hướng xuống.
\item Phương trình vận tốc: $v=v_0+at=10t$.
\item Vận tốc lúc chạm đất $v = 50$ m/s $\Rightarrow$ thời gian vật rơi là: $t=\dfrac{50}{10}=5\ s$. 
\item Tọa độ vật: $y=\dfrac{gt^2}{2}=5t^2$. 
\item Vật được thả từ độ cao: $h=5t^2=5.5^2=125\ m$. 
\item Quãng đường vật đi được sau 3 s là: $s=5.3^2=45\ m$.
\item Độ cao của vật sau khi đi được 3s là: $h'=h-s=125-45=80\ m$.
\end{itemize}	
}
\haidong{7} \end{vidu}
\begin{vidu}
Thời gian để một quả bóng rơi từ đỉnh tháp có chiều cao $h$ xuống đất là $t$.  Bỏ qua sức cản không khí. Tại thời điểm $\dfrac{t}{3}$ sau khi quả bóng rơi, quả bóng cách mặt đất một khoảng
\choice
{$\dfrac{h}{9}$}
{$\dfrac{7 h}{9}$}
{\True $\dfrac{8 h}{9}$}
{$\dfrac{17 h}{18}$}
\haidong{5}
\loigiai{
Quãng đường vật rơi trong $t s: h=\dfrac{1}{2} g \cdot t^2$

Quãng đường vật rơi trong $\dfrac{t}{3} s$ đầu tiên: $h_1=\dfrac{1}{2} g \cdot\left(\dfrac{t}{3}\right)^2=\dfrac{h}{9}$ 

Độ cao của vật khi đó: $h^{\prime}=h-h_1=\dfrac{8 h}{9}$

}
\haidong{7} \end{vidu}
\loai{Thay đổi thông số}
\begin{vidu} %[1P1-7.2B][Loai1]
Bỏ qua sức cản không khí. Thả hòn đá từ độ cao $h$ xuống đất. Hòn đá rơi trong $1\ s$. Nếu thả hòn đá đó từ độ cao $4h$ xuống đất thì hòn đá sẽ rơi trong bao lâu?
\choice
{$4\ s$}
{\True $2\ s$}
{$\sqrt2\ s$}
{$3\ s$}
\haidong{4}
\loigiai{		
\begin{itemize}
\item Ta có lần rơi thứ nhất: $t_1 = \sqrt {\dfrac{2h}{g}} = 1\ s$.
\item Lần rơi thứ hai: $t_2 = \sqrt {\dfrac{8h}{g}} = 2\sqrt{\dfrac{2h}{g}} = 2t_1 = 2\ s$.
\end{itemize}	
}
\haidong{7} \end{vidu}
\begin{vidu} %[1P1B7-1][Loai1]
Trong khi rơi tự do, vật thứ nhất rơi mất một khoảng thời gian dài gấp đôi vật thứ hai. So sánh quãng đường đi được của vật thứ nhất và vật thứ hai.
\choice
{$h_1 = \dfrac{1}{2}h_2$}
{$h_1 = 2h_2$}
{$h_1 = 3h_2$}
{\True $h_1 = 4h_2$}
\haidong{3}
\loigiai{
\begin{itemize}
\item Với vật thứ nhất ta có: $h_1 = \dfrac{1}{2}gt_1^2$.
\item Với vật thứ hai: $t_2 = \dfrac{1}{2}t_1$.
\item Quãng đường đi được của vật thứ hai là: $h_2 = \dfrac{1}{2}gt_2^2 = \dfrac{1}{2}g\left( {\dfrac{1}{2}t_1} \right)^2 \Rightarrow h_1 = 4h_2$.
\end{itemize}	
}
\haidong{7} \end{vidu}
\loai{Tính tốc độ trung bình}
\begin{vidu} %[1P1K7-2][Loai1]
Lấy $g = 10\  m/s^2$. Một vật rơi tự do thì tốc độ trung bình của vật trong quãng đường $20\ m$ đầu tiên là
\choice
{\True $10\  m/s$}
{$20\  m/s$}
{$10\sqrt 2\  m/s$}
{$20\sqrt 2\  m/s$}
\haidong{5}
\loigiai{	      
\begin{itemize}
\item Tốc độ trung bình: $v_{tb} = \dfrac{s}{t} = \dfrac{h}{t}$.
\item Mà ta lại có: $t = \sqrt {\dfrac{2h}{g}} = \sqrt{\dfrac{2.20}{10}} = 2\ s.$
\item $v_{tb} = \dfrac{{20}}{2} = 10\ m/s.$
\end{itemize}	
}
\haidong{7} \end{vidu}
\loai{Hai vật rơi khác thời điểm}
\begin{vidu}%[1P1K7-2][Loai1]
Hai viên bi sắt được thả rơi cùng độ cao cách nhau một khoảng thời gian $0,5\ s$. Lấy $g = 10\ m/s^2$. Khoảng cách giữa hai viên bi sau khi viên thứ nhất rơi được $1,5\ s$ là
\choice
{\True $6,25\ m$}
{$12,5\ m$}
{$5,0\ m$}
{$2,5\ m$}
\haidong{6}
\loigiai{
\begin{itemize}
\item Chọn gốc tọa độ là vị trí thả hai viên bi, chiều dương hướng xuống.
\item Phương trình chuyển động của 2 viên bi là: $y_1=\dfrac{gt_1^2}{2}\ m;\ y_2=\dfrac{gt_2^2}{2}\ m$.
\item Sau khi viên bi thứ nhất rơi được 1,5 s thì viên bi thứ hai rơi được 1 s.
\item Ta có: $y_1=\dfrac{10.1,5^2}{2}=11,25\ m;\ y_2=\dfrac{10.1^2}{2}=5\ m$.
\item Khoảng cách giữa hai viên bi là: $y=y_1-y_2=11,25-5=6,25\ m$.
\end{itemize}	
}
\haidong{7} \end{vidu}
\begin{vidu}%[1P1-7.2K][Loai1]
Bỏ qua sức cản không khí. Lấy $g = 10\ m/s^2$. Từ một đỉnh tháp, bạn Hưng buông rơi tự do một vật. Một giây sau ở tầng tháp thấp hơn $10\ m$, bạn Tùng buông rơi tự do vật thứ $2$. Kể từ khi thả vật $1$ thì hai vật gặp nhau sau
\choice
{$1,2$ giây}
{\True $1,5$ giây}
{$2,5$ giây}
{$3,5$ giây}
\haidong{8}
\loigiai{
\immini[thm]{
\begin{itemize}
\item Quãng đường đi được của vật 1: 	$s_1=\dfrac{1}{2}gt^2\Rightarrow s=5t^2$.
\item Quãng đường đi được của vật 2: 	$s_2=\dfrac{1}{2}g{\left(t-1\right)}^2\Rightarrow s_2=5{\left(t-1\right)}^2$.
\item Khi hai vật gặp nhau: $s_1-s_2 =10 \Rightarrow  t = 1,5$ s.
\end{itemize}
}{\begin{tikzpicture}[scale=1,thick,>=stealth']
\tkzDefPoints{0/0/o,0/-2/a, 0/-5/y, 0.2/0/s}
\tkzDrawPoints[size=2](o,a)
\tkzDrawSegments[->](o,y)
\tkzLabelPoint[left](o){$O$}
\tkzLabelPoint[left](a){$A$}
\tkzLabelPoint[right=5pt](o){$t=0$}
\tkzLabelPoint[right=5pt](a){$t=1\ s$}
\begin{scope}[/pgf/decoration/raise=5pt]
\draw [decorate,decoration={brace,mirror,amplitude=5pt},xshift=0pt,yshift=-4pt]
(a) -- (o) node [black,midway,xshift=25pt] 
{$10\ m$};
\end{scope}
\end{tikzpicture}
}
}
\haidong{7} \end{vidu}
\loai{Quãng đường vật rơi trong giây thứ n}
\begin{vidu}%[1P1K7-3][Loai1]
Một vật rơi tự do không vận tốc ban đầu từ độ cao $180\ m$ xuống.  Cho $g = 10\  m/s^2$. Quãng đường vật đi được trong giây cuối cùng là
\choice
{$30\ m$}
{$45\ m$}
{\True $55\ m$}
{$125\ m$}	
\haidong{5}
\loigiai{
\begin{itemize}
\item Chọn gốc tọa độ tại vị trí vật rơi, chiều dương hướng xuống.
\item 	Tọa độ của vật: $y=\dfrac{gt^2}{2}=5t^2$
\item 	Khi chạm đất thì  $y=180=5t^2 
\Rightarrow$ thời gian vật rơi $t = 6$ s.
\item 	Quãng đường vật đi được trong 5 s đầu tiên là: $s=5t^2=5.5^2=125\ m$.
\item 	Quãng đường vật đi được trong giây cuối cùng là: $s'=h-s=180-125=55\ m$.
\end{itemize}	
}
\haidong{7} \end{vidu}
\begin{vidu} %[1P1K7-3][Loai1] %BAI 4+5 - HOC MAI
Thả một hòn đá từ một độ cao $h$. Cho $g = 10\  m/s^2$. Tính h biết rằng trong giây cuối cùng hòn đá đi được $25\ m$. 
\choice
{\True $45\ m$}
{$40\ m$}
{$35\ m$}
{$50\ m$}
\haidong{5}
\loigiai{
\begin{itemize}
\item Ta có $h = 0,5gt^2;\ h_1 = 0,5.g(t - 1)^2$
\item Quãng đường đi được trong giây cuối là $\Delta h = h - h_1 = 0,5.g(2t - 1) = 25 \Rightarrow t = 3$ s hay độ cao $h = 45$ m.
\end{itemize}	
}
\haidong{7} \end{vidu}
\begin{vidu}%[1P1K7-3][Loai1]
Một vật $A$ được thả rơi từ độ cao $45\ m$ xuống mặt đất. Lấy $g = 10\ m/s^2$. Quãng đường vật rơi được trong $2$ giây cuối cùng là	
\choice
{\True $40\ m$}
{$35\ m$}
{$30\ m$}
{$25\ m$}
\haidong{6}
\loigiai{
\begin{itemize}
\item Chọn gốc tọa độ tại vị trí vật rơi, chiều dương hướng xuống dưới.
\item Ta có: $y=\dfrac{gt^2}{2}=5t^2\ m$
\item Thời gian vật rơi là nghiệm phương trình: $5t^2=45 \Rightarrow t=3\ s$
\item Quãng đường vật rơi trong 2 s cuối cùng là: $s=y_3-y_1=45-5.1^2=40\ m$.

\end{itemize}	
}
\haidong{7} \end{vidu}
\begin{vidu}
Một vật rơi tự do trong thời gian $t$. Trong nửa thời gian đầu vật rơi được quãng đường $s_1$, trong nửa thời gian sau vật rơi được quãng đường $s_2$. Hệ thức nào sau đây là đúng?
\choice
{$s_2=s_1$}
{$s_2=2 s_1$}
{\True $s_2=3 s_1$}
{$s_2=4 s_1$}
\haidong{6}
\loigiai{
Quãng đường vật rơi trong nửa thời gian đầu $s_1=\dfrac{1}{2} g t_1^2$

Quãng đường vật rơi trong tổng thời gian $s=\dfrac{1}{2} g t^2=\dfrac{1}{2} g \cdot\left(2 t_1\right)^2=4 s_1$

$\rightarrow$ Quãng đường vật rơi trong nửa thời gian sau là $s_2=s-s_1=3 s_1$

}
\haidong{7} \end{vidu}

\loai{Rơi tự do và chuyển động khác}
\begin{vidu}%[1P1-7.2K][Loai1]
Để biết độ sâu của một cái giếng đã hết nước, người ta thả một hòn đá từ miệng giếng và đo thời gian từ lúc thả đến lúc nghe thấy tiếng vọng của hòn đá khi chạm đáy. Giả sử người ta đo được thời gian là $2,06\ s$. Lấy gia tốc trọng trường $g = 10\ m/s^2$ và vận tốc âm trong không khí là $340\ m/s$. Coi âm truyền theo một phương nào đó là thẳng đều. Độ sâu của giếng ước lượng là
\choice
{$10\ m$}
{$15\ m$}
{$30\ m$}
{\True $20\ m$}
\haidong{4}
\loigiai{
\begin{itemize}
\item Thời gian rơi tự do $t_1$ của hòn đá: $h=\dfrac{1}{2}gt_1^2\Rightarrow t_1=\sqrt{\dfrac{2h}{g}}=\sqrt{\dfrac{h}{5}}$.
\item Thời gian truyền âm $t_2$: $t_2=\dfrac{h}{v_m}=\dfrac{h}{340}$.
\item Theo đề ta có: $t_1+t_2=2,06\Leftrightarrow \dfrac{h}{340}+\sqrt{\dfrac{h}{5}}=2,06\Rightarrow h\approx 20\ m$.
\end{itemize}
}
\haidong{7} \end{vidu}
